\section*{الفصل \ref{ch:g8:winds-global-circulation} --- الرياح والدورة العامة للغلاف الجوي}

\begin{solution}{exo:g8:winds-global-circulation:1}
(أ) الريح التجارية الشمالية الشرقية، من الشمال الشرقي. (ب) \omterm{def:g8:winds-global-circulation:global-winds}{الرياح الغربية}،
من الشمال الغربي. (ج) \omterm{def:g8:winds-global-circulation:global-winds}{الرياح الشرقية القطبية}، من الشرق (من الشمال
الشرقي). (د) الريح التجارية الجنوبية الشرقية، من الجنوب الشرقي.
\end{solution}

\begin{solution}{exo:g8:winds-global-circulation:2}
(أ) الهواء الدافئ يصعد والهواء البارد يهبط. (ب) منخفضًا. (ج) من الضغط العالي
نحو الضغط المنخفض. (د) إلى اليمين.
\end{solution}

\begin{solution}{exo:g8:winds-global-circulation:3}
(أ) خطأ: \omterm{def:g8:winds-global-circulation:breezes}{نسيم البحر} يهب من \omterm{def:g2:water-on-earth:ocean}{البحر} إلى اليابسة (في النهار).
(ب) صواب. (ج) خطأ: فهي تهب من الشرق (من الشمال الشرقي أو من الجنوب
الشرقي). (د) صواب.
\end{solution}

\begin{solution}{exo:g8:winds-global-circulation:4}
يُظهر الرسم حلقة: على طول الأرض، سهم من \omterm{def:g2:water-on-earth:ocean}{البحر} إلى اليابسة؛ وفوق اليابسة،
سهم صاعد؛ وفي الأعلى، سهم من اليابسة عائد فوق \omterm{def:g2:water-on-earth:ocean}{البحر}؛ وفوق \omterm{def:g2:water-on-earth:ocean}{البحر}، سهم
هابط. ويوضع الحرف ع على الأرض تحت الهواء الهابط (فوق \omterm{def:g2:water-on-earth:ocean}{البحر})، والحرف م تحت
الهواء الصاعد (فوق اليابسة).
\end{solution}

\begin{solution}{exo:g8:winds-global-circulation:5}
(أ) 1.41 ضعفًا من الأرض. (ب) 50~\%. (ج) 26~\%. (د) قرب \omterm{def:g3:continents-and-oceans:poles}{القطبين}،
يصل ضوء الشمس مائلًا وينتشر على مساحة كبيرة، فيتلقى كل متر مربع قليلًا
من الطاقة.
\end{solution}

\begin{solution}{exo:g8:winds-global-circulation:6}
(أ) \omterm{def:g8:winds-global-circulation:global-winds}{الرياح الشرقية القطبية} (من \ang{70} إلى \ang{60} شمالًا)، \omterm{def:g8:winds-global-circulation:global-winds}{فالرياح الغربية} (من \ang{60} إلى
\ang{30} شمالًا)، \omterm{def:g8:winds-global-circulation:global-winds}{فالرياح التجارية} الشمالية الشرقية (من \ang{30} شمالًا إلى \omterm{def:g3:continents-and-oceans:equator}{خط الاستواء})،
\omterm{def:g8:winds-global-circulation:global-winds}{فالرياح التجارية} الجنوبية الشرقية (من \omterm{def:g3:continents-and-oceans:equator}{خط الاستواء} إلى \ang{30} جنوبًا)، \omterm{def:g8:winds-global-circulation:global-winds}{فالرياح الغربية} (من \ang{30} إلى
\ang{60} جنوبًا)، \omterm{def:g8:winds-global-circulation:global-winds}{فالرياح الشرقية القطبية} (من \ang{60} إلى \ang{70} جنوبًا). (ب) قرب \omterm{def:g3:continents-and-oceans:equator}{خط
الاستواء}، في منطقة الركود الاستوائي. (ج) قرب \ang{30} شمالًا و\ang{30} جنوبًا، تحت
الهواء الهابط.
\end{solution}

\begin{solution}{exo:g8:winds-global-circulation:7}
(أ) اليابسة ($31 > 22$): تهب الريح من \omterm{def:g2:water-on-earth:ocean}{البحر} إلى اليابسة، إنه \omterm{def:g8:winds-global-circulation:breezes}{نسيم
البحر}. (ب) \omterm{def:g2:water-on-earth:ocean}{البحر} ($21 > 17$): تهب الريح من اليابسة إلى \omterm{def:g2:water-on-earth:ocean}{البحر}، إنه \omterm{def:g8:winds-global-circulation:breezes}{نسيم
البر}. (ج) اليابسة: $31 - 17 = 14$ درجة؛ \omterm{def:g2:water-on-earth:ocean}{والبحر}: $22 -
21 = 1$ درجة. (د) أشعة الشمس لا تسخّن إلا طبقة رقيقة من الأرض،
تسخن وتبرد بسرعة؛ أما \omterm{def:g2:water-on-earth:ocean}{البحر} فينشر حرارته في طبقة سميكة من الماء، فلا تكاد
درجة حرارته تتغير بين النهار والليل.
\end{solution}

\begin{solution}{exo:g8:winds-global-circulation:8}
(أ) من الشمال. (ب) إلى اليمين. (ج) من الشمال الشرقي (إنها الريح التجارية
الشمالية الشرقية). (د) لولا الدوران لهب من الجنوب؛ والدوران يحرفه إلى
اليسار؛ فيهب من الجنوب الشرقي.
\end{solution}

\begin{solution}{exo:g8:winds-global-circulation:9}
قرب \omterm{def:g3:continents-and-oceans:equator}{خط الاستواء}، تسخّن الشمس الأرض \omterm{def:g2:water-on-earth:ocean}{والبحر} بشدة، فالهواء \emph{يسخن}؛
والهواء الدافئ الرطب \emph{يصعد}؛ وكلما صعد \emph{يبرد}،
فيتحول بخار الماء الذي فيه إلى \emph{\omterm{def:g4:clouds-rain-snow:cloud}{سحب}}، تعطي زخات وعواصف رعدية. وقرب
\ang{25} شمالًا، الهواء الذي صعد عند \omterm{def:g3:continents-and-oceans:equator}{خط الاستواء} \emph{يهبط}؛ وتحته يكون
الضغط عند الأرض \emph{عاليًا}؛ والهواء الهابط يسخن، وتتبدد سحبه
وتبقى المنطقة \emph{جافة}: إنها \omterm{def:g5:weather-and-climate:desert}{صحراء}.
\end{solution}

\begin{solution}{exo:g8:winds-global-circulation:10}
(أ) \omterm{def:g8:winds-global-circulation:global-winds}{الرياح الغربية}. لا: فهي تهب من الغرب، معاكسة لسفينة تبحر غربًا. (ب)
لتبلغ حزام \omterm{def:g8:winds-global-circulation:global-winds}{الرياح التجارية}، جنوب نحو \ang{30} شمالًا. (ج) الريح
التجارية الشمالية الشرقية، بين \omterm{def:g3:continents-and-oceans:equator}{خط الاستواء} ونحو \ang{30} شمالًا. (د) \omterm{def:g8:winds-global-circulation:global-winds}{الرياح
الغربية}، شمال نحو \ang{30} شمالًا (قرب \ang{40} شمالًا، دائرة عرض ميناء
الانطلاق).
\end{solution}

\begin{solution}{exo:g8:winds-global-circulation:11}
(أ) $1 \div 0.5 = 2$: والجدول يعطي 2. (ب) $1 \div 0.866 \approx
1.1547$، أي \textbf{1.15}. (ج) عند \ang{30}: $1000 \div 1.15 \approx
\qty{870}{W}$ على كل متر مربع (وبالضبط $1000 \times 0.866 =
\qty{866}{W}$). وعند \ang{60}: $1000 \div 2 = \qty{500}{W}$.
\end{solution}

\begin{solution}{exo:g8:winds-global-circulation:12}
(أ) ثلاث حلقات: يصعد الهواء عند \omterm{def:g3:continents-and-oceans:equator}{خط الاستواء} ويهبط قرب \ang{30}؛
ويصعد قرب \ang{60} ويهبط قرب \ang{30}؛ ويصعد قرب \ang{60} ويهبط
عند \omterm{def:g3:continents-and-oceans:poles}{القطب}. (ب) م عند \omterm{def:g3:continents-and-oceans:equator}{خط الاستواء} وقرب \ang{60}؛ وع قرب \ang{30} وعند
\omterm{def:g3:continents-and-oceans:poles}{القطب}. (ج) من \ang{30} إلى \omterm{def:g3:continents-and-oceans:equator}{خط الاستواء}، يتجه الهواء السطحي جنوبًا
وينحرف إلى اليمين: الريح التجارية الشمالية الشرقية. ومن \ang{30} إلى
\ang{60}، يتجه شمالًا وينحرف إلى اليمين: فيأتي من الجنوب
الغربي. ومن \omterm{def:g3:continents-and-oceans:poles}{القطب} إلى \ang{60}، يتجه جنوبًا، منحرفًا إلى
اليمين: الريح الشرقية القطبية. (د) في الخلية الوسطى يتدفق الهواء السطحي
نحو \omterm{def:g3:continents-and-oceans:poles}{القطب}، أي شمالًا؛ ومنحرفًا إلى اليمين، يتجه نحو
الشمال الشرقي: تلك هي \omterm{def:g8:winds-global-circulation:global-winds}{الرياح الغربية}.
\end{solution}

\begin{solution}{pb:g8:winds-global-circulation:1}
\textbf{1.} $59.6 - 23.5 = 36.1$ درجة من \omterm{def:g6:latitude-longitude:longitude}{خطوط الطول}.

\textbf{2.} $36.1 \times 108 \approx 3899$: نحو \qty{3900}{km}.

\textbf{3.} نعم: $3900 - 3896 = 4$~كيلومترات، وهو فرق ضئيل.

\textbf{4.} نعم: فكلاهما يقع قرب \ang{14} شمالًا، بين \omterm{def:g3:continents-and-oceans:equator}{خط الاستواء} ونحو
\ang{30} شمالًا، وهو حزام الريح التجارية الشمالية الشرقية.

\textbf{5.} من الشمال الشرقي (فالأسهم تشير نحو
الجنوب الغربي).

\textbf{6.} إلى الأمام، وجانبيًا قليلًا: فالريح تأتي من خلف
السفينة، ومن يمينها قليلًا، وتدفعها نحو الغرب.

\textbf{7.} لأنها جزء من \omterm{def:g8:winds-global-circulation:circulation}{الدورة العامة}: فالهواء لا يزال يهبط قرب
\ang{30} شمالًا ويصعد قرب \omterm{def:g3:continents-and-oceans:equator}{خط الاستواء}، فيظل الهواء السطحي يتدفق
نحو \omterm{def:g3:continents-and-oceans:equator}{خط الاستواء}، ودوران الأرض يحرفه دائمًا إلى
اليمين.

\textbf{8.} لأن الريح التجارية ستهب معاكسة للسفينة. فينبغي للربان أن
يبحر شمالًا، إلى ما بعد نحو \ang{30} شمالًا، ليجد \omterm{def:g8:winds-global-circulation:global-winds}{الرياح الغربية}.

\textbf{9.} $7 \times 24 = \qty{168}{km}$ في اليوم.

\textbf{10.} $3 \times 168 = \qty{504}{km}$.

\textbf{11.} $3896 \div 168 \approx 23.2$ يومًا.

\textbf{12.} نحو \textbf{23 يومًا}: أقل من شهر.
\end{solution}
